\documentclass[12pt]{article}
\usepackage[utf8]{inputenc}
\usepackage[margin=1in]{geometry}
\usepackage{amsmath, amssymb}
\usepackage{graphicx}
\usepackage{hyperref}

\title{Project Report Sample}
\author{Author Name}
\date{\today}

\begin{document}

\maketitle

\begin{abstract}
This document provides a sample structure for a project report written in \LaTeX. It serves as a template for academic or technical projects.
\end{abstract}

\tableofcontents

\section{Introduction}
Introduce the motivation for the project, summarize the problem statement, and outline the objectives of the work.

\section{Background}
Provide relevant literature, theoretical background, or technical context required to understand the project.

\section{Methods}
Describe the methodology used in the project. Detail the models, algorithms, or processes applied.

\subsection{Data Collection}
Explain how data was collected or sourced for the project.

\subsection{Implementation Details}
Discuss key implementation details. Include code snippets, figures, or flowcharts where helpful.

\section{Results}
Present the main findings. Use figures, tables, and graphs as needed.

\begin{figure}[h!]
    \centering
    \includegraphics[width=0.6\linewidth]{figures/sample-figure.png}
    \caption{Sample figure.}
    \label{fig:sample}
\end{figure}

\section{Discussion}
Interpret the results, discuss implications, limitations, and compare with related work.

\section{Conclusion}
Summarize the achievements of the project and suggest avenues for future work.

\section*{References}
% Add your references here, for example:
% \bibliographystyle{plain}
% \bibliography{references}

\end{document}